\section{算法细节与工程权衡：On-policy、更新强度与熵正则}

\subsection{On-policy 与 Off-policy 的核心差异}
讲者用直观方式解释了 on-policy：模型应优先从 ``自己当前策略生成的轨迹'' 中获得反馈。因为这些错误最贴近当前能力边界，学习效率更高。对 Agent 训练来说，这通常能更好地维持策略一致性和探索质量。

\begin{knowledgebox}{工程上如何理解 on-policy 优势}
on-policy 的价值不是理论口号，而是减少 ``反馈错位''：你给模型的奖励，最好对应模型当前真实会走的轨迹，而不是别的策略分布。这样更新方向更一致，避免学到不可复现行为。
\end{knowledgebox}

\subsection{更新强度控制：clipping 与解耦阈值}
课程讨论了 RL 更新中的 clipping 调整：通过上/下界非对称设置，鼓励低概率 token 的合理探索，缓解过早收敛。讲者也坦诚这类技巧在当下更像 ``有效工程 hack''，未必是长期最优理论形式。

\begin{importantbox}{工程结论}
在可验证 Agent 训练里，\textbf{稳定更新} 与 \textbf{保持探索} 同等重要。只追求稳定会导致保守退化，只追求探索会导致训练发散。clipping、熵正则等机制本质上是在做这两个目标的平衡。
\end{importantbox}

\subsection{熵正则：显式维持可探索空间}
当观测到熵持续下降时，常见做法是加入熵相关损失项，将生成熵维持在合理区间。讲座中提到社区已有多种实现方案，效果依赖任务、模型与 verifier 质量，尚不存在一次性通吃的方法。

\begin{warningbox}{不要把技巧当定律}
当前许多有效技巧是 ``经验可行''，并不代表理论已收敛。团队在迁移方法时应保留审计和消融流程，避免把某个数据集上的增益误判为普适规律。
\end{warningbox}

\subsection{本章小结}
算法层面最关键的不是某个单一公式，而是 ``反馈一致性 + 更新稳定性 + 探索维持'' 三者协同。on-policy、clipping 调整、熵正则都是围绕这三点展开。

